\section{Local search and anchored comparison clusters}
\label{loc:section}

We first prove the two local statements used by both branches of the
algorithm.  The geometric arguments apply to any finite metric.  The bounded
integral distance assumption is needed only for the polynomial running time
of exact local descent.
Exchanges of bounded size are standard in local-search analysis
\cite[Sections~3.3--3.4]{Arya2004}.  Our weighted comparison adapts the
nearest-facility routing proof of Gupta and Tangwongsan
\cite{GuptaTangwongsan2008} while protecting the specified proxies.
The anchoring argument adapts the disjoint-region and incoming-client
analysis of \cite{CohenAddadSchwiegelshohn2017} and
\cite[Section~4.1.1]{CGLSS2026}.  We prove the particular scale
hypothesis and error bound needed here, for a comparison solution that
need not be optimal.

\subsection{Local optimality and weighted exchanges}

Fix an integer $r$ with $1\le r\le |F|$.  A set $S\subseteq F$ of size $r$
is a \emph{width-$5$ local optimum} if no exchange removing at most five
facilities and adding at most five facilities produces a set of size $r$
with strictly smaller cost.  We also use exchanges whose resulting set is
nonempty and has fewer than $r$ facilities: any such set can be padded to
size $r$ by restoring removed facilities, without increasing its cost or
the width of the exchange.

\begin{lemma}[Exact local descent]
\label{loc:descent}
Suppose all distances are integers at most $B$, where $B$ is bounded by a
fixed polynomial in $N=n+|F|$.  From any nonempty initial set of $r$
facilities, a width-$5$ local optimum of size $r$ can be found in polynomial
time, with no increase in cost.
\end{lemma}

\begin{proof}
Enumerate all sets of at most five facilities to remove and all sets of at
most five facilities to add, checking those exchanges that leave exactly
$r$ facilities.  There are $O(N^{10})$ pairs, and each resulting cost can
be evaluated in polynomial time.  Perform any strict improvement and
repeat.  Every cost is an integer between zero and $nB$, so at most $nB$
strict improvements are performed.

For completeness, consider an exchange with removal set $R\subseteq S$,
addition set $A\subseteq F$, $|R|,|A|\le5$, and nonempty result
$T=(S\setminus R)\cup A$ of size at most $r$.  All facilities in $S\setminus T$
belong to $R$, and
$|S\setminus T|\ge r-|T|$.  Restoring any $r-|T|$ of them gives a set of
size $r$, still differing from $S$ by at most five removals and five
additions, with cost at most $\operatorname{cost}(T)$.  Thus the final local
optimum also rules out every improving exchange of this form.
\end{proof}

Let $O\subseteq F$ be a comparison set of size $h$, where
$1\le h\le |S|$, and put $P=\operatorname{cost}(O)>0$.  Fix nearest-facility
assignments to both $O$ and $S$, using arbitrary fixed rules for ties.
For $i\in O$, write $C_i$ for its assigned client set and define
\[
 n_i=|C_i|,\qquad
 d_p^*=d(p,O),\qquad s_p=d(p,S),\qquad
 P_i=\sum_{p\in C_i}d_p^*.
\]
Write $s(p)\in S$ for the assigned local facility, so
$d(p,s(p))=s_p$.  Empty sets $C_i$ are permitted.  Given a partition
$O=O_g\mathbin{\dot\cup}O_b$, define
\[
 P_g=\sum_{i\in O_g}P_i,\quad P_b=\sum_{i\in O_b}P_i,\quad
 S_g=\sum_{i\in O_g}\sum_{p\in C_i}s_p,\quad
 S_b=\sum_{i\in O_b}\sum_{p\in C_i}s_p.
\]
For each good center $i\in O_g$, suppose a proxy $\phi_i\in S$ has been
specified, with these proxies pairwise distinct.  Set
\[
 D_g=\sum_{i\in O_g}\sum_{p\in C_i}d(p,\phi_i).
\]

\begin{lemma}[Weighted swap bound]
\label{loc:swap}
Under the preceding hypotheses, if $S$ is a width-$5$ local optimum and
$u=6/5$, then
\begin{equation}
 S_b+uS_g\le (1+2u)P_b+uD_g.
 \label{loc:strong-swap}
\end{equation}
Consequently,
\begin{equation}
 \operatorname{cost}(S)\le uD_g+(1+2u)P_b.
 \label{loc:swap-bound}
\end{equation}
In particular, for any comparison set $O$ of size at most $|S|$,
$\operatorname{cost}(S)\le(17/5)P\le5P$.
\end{lemma}

\begin{proof}
Protect the proxy set $\Phi=\{\phi_i:i\in O_g\}$, meaning that no exchange
below removes any member of $\Phi$.  Write $g=|O_g|$, $m=|O_b|$, and
$q=|S\setminus\Phi|$.  Distinctness of the proxies and $h\le|S|$ imply
\[
 q=|S|-g\ge h-g=m.
\]
If $m=0$, every client can use its proxy, so $S_g\le D_g$ and both
inequalities follow.  Assume henceforth that $m>0$; in particular an
unprotected facility exists.

For every $i\in O_b$, choose a nearest facility $\nu(i)\in S$ to $i$.
Retain this choice for the routing argument.  Separately define an image
$\psi(i)\in S\setminus\Phi$: if $\nu(i)$ is unprotected, take
$\psi(i)=\nu(i)$; if $\nu(i)$ is protected, take any unprotected facility
as $\psi(i)$.  Thus the image used to organize exchanges need not be the
nearest facility used for routing.

Let $a$ be the number of nonempty image classes of $\psi$.  A class of
size $t$ receives its image facility together with $t-1$ auxiliary
unprotected facilities that are not images of any bad center.  These sets
of facilities can be chosen disjointly across classes: the total number
of auxiliary facilities needed is $m-a$, whereas the number available is
$q-a\ge m-a$.  Facilities in these allocated sets will be called slots.

We give a nonnegative weighting of feasible exchanges.  For an image
class $I\subseteq O_b$ of size $t\le5$, make the single exchange of weight
one that removes its $t$ slots and adds every center in $I$.  For a class
of size $t>5$, for each $i\in I$ and each of the $t-1$ auxiliary slots,
make the exchange that removes that slot and adds $i$, with weight
$1/(t-1)$.  These are exchanges of width at most five.  Their results are
nonempty and have size at most $|S|$, even when an added facility already
belongs to $S$: they are interpreted as actual set unions, and
Lemma~\ref{loc:descent}'s padding observation applies.  Local optimality
therefore gives a nonnegative cost change for every one of them.

Every bad center is added in exchanges of total weight exactly one.  A
slot allocated to a class of size at most five has removal weight one.
In a class of size $t>5$, the image slot is never removed and each
auxiliary slot has removal weight
\[
 \frac{t}{t-1}\le\frac65=u.
\]
Unallocated facilities are never removed.  Thus every facility in $S$
has total removal weight at most $u$.

We also need a preservation property.  In every exchange that does not
add a particular bad center $i$, its remembered nearest facility
$\nu(i)$ survives.  If $\nu(i)$ is protected, this is part of the
construction.  Otherwise $\nu(i)=\psi(i)$ is the image slot of $i$'s
class.  A small-class exchange that removes this slot adds all centers
of the class, including $i$; a large-class exchange never removes it.
Other classes have disjoint slots.  This proves the property, including
when $O\cap S\ne\varnothing$.

Consider a client $p\in C_i$ with $i\in O_b$.  In every exchange that
adds $i$, route $p$ to $i$, at cost change $d_p^*-s_p$.  In any other
exchange retain its old assignment if possible.  If its old facility
is removed, use the preserved facility $\nu(i)$, for which
\begin{align*}
 d(p,\nu(i))
 &\le d_p^*+d(i,\nu(i))\\
 &\le d_p^*+d(i,s(p))
 \le 2d_p^*+s_p.
\end{align*}
The weighted sum of these upper bounds for $p$ is therefore at most
\[
 d_p^*-s_p+2u d_p^*.
\]
Here the first term uses total addition weight one.  For the other
exchanges we upper-bound their removal weight by the full removal
weight of $s(p)$, which is at most $u$; the extra charge $2d_p^*$ is
nonnegative.  Removing and immediately adding back the same facility
only reduces the need for this extra charge.

For a client $p\in C_i$ with $i\in O_g$, keep the old assignment unless
its serving facility is removed, in which case use the protected
proxy $\phi_i$.  Nearest assignment in $S$ gives
$d(p,\phi_i)-s_p\ge0$, so its weighted cost change is at most
$u(d(p,\phi_i)-s_p)$.

Each described routing is feasible for its exchange.  Summing the
upper bounds on cost change, and using nonnegativity of every actual
cost change, yields
\[
 0\le P_b-S_b+2uP_b+u(D_g-S_g).
\]
This is \eqref{loc:strong-swap}.  Since $u\ge1$, it implies
\eqref{loc:swap-bound}.  Declaring every center bad gives the final
assertion.
\end{proof}

\subsection{Anchoring at any fixed accuracy}

For $h\ge2$ define the separation of a comparison center by
\[
 \operatorname{sep}_i
 =\min_{j\in O\setminus\{i\}}d(i,j),
 \qquad \ell=|S|-h\ge0.
\]
In a metric, $\operatorname{sep}_i>0$.  The next theorem is existential:
the algorithm will not need to find its partition or proxies.

\begin{theorem}[Anchoring]
\label{loc:anchoring}
Let $S$ be a width-$5$ local optimum, and let $O$ be a comparison set
with $2\le h=|O|\le|S|$ and $P=\operatorname{cost}(O)>0$.  Suppose $b>0$
satisfies
\begin{equation}
 b\left(\ell+
   \bigl|\{i\in O:n_i\operatorname{sep}_i<b\}\bigr|\right)
 \le10P.
 \label{loc:scale-condition}
\end{equation}
For every fixed $\mu\in(0,1]$, there is a partition
$O=O_g\mathbin{\dot\cup}O_b$ and distinct proxies $\phi_i\in S$ for
$i\in O_g$ such that, with $m=|O_b|$,
\begin{equation}
 \sum_{i\in O_g}\sum_{p\in C_i}d(p,\phi_i)
 \le P_g+\mu P,
 \qquad m\le C_\mu\frac{P}{b},
 \label{loc:anchoring-conclusion}
\end{equation}
where the finite constant
\begin{equation}
 J_\mu=\left\lceil\frac{40}{\mu}\right\rceil,
 \qquad
 C_\mu=10+\frac{600}{\mu}\,2^{J_\mu-1}
 \label{loc:anchoring-constant}
\end{equation}
depends only on $\mu$.
\end{theorem}

\begin{proof}
Set
\[
 V_i=\sum_{p\in C_i}(s_p+d_p^*),\qquad
 q_0=\bigl|\{i\in O:n_i\operatorname{sep}_i<b\}\bigr|.
\]
The all-bad case of Lemma~\ref{loc:swap} gives
\begin{equation}
 \sum_{i\in O}V_i=\operatorname{cost}(S)+P\le6P.
 \label{loc:V-bound}
\end{equation}
Call a center \emph{short} if $n_i\operatorname{sep}_i<b$.  In particular,
every empty comparison cluster is short.

\paragraph{Selecting a threshold.}
Let $t_a=(\mu/100)2^{-a}$ for $0\le a<J_\mu$.  At threshold $t_a$, declare
$i$ good precisely when it is not short and $V_i\le t_ab$, and let
$m_a$ be the resulting number of bad centers.  For a fixed nonshort
center $i$, the geometric tail of those values $t_ab$ that are strictly
less than $V_i$ has sum at most $2V_i$.  Consequently,
\begin{align}
 \sum_{a=0}^{J_\mu-1}t_ab(m_a+\ell)
 &\le 2(\mu/100)b(q_0+\ell)+2\sum_{i\in O}V_i\notag\\
 &\le (12+\mu/5)P.
 \label{loc:threshold-average}
\end{align}
The first term accounts for short centers and surplus facilities at
every threshold; the second accounts for all other bad centers.
Since $J_\mu\ge40/\mu$ and $\mu\le1$, the average of the left-hand
summands is at most
\[
 \frac{\mu}{40}(12+\mu/5)P
 \le\frac{61}{200}\mu P<\frac{\mu P}{2}.
\]
Choose a threshold $t=t_a$ attaining at most this average, and use its
good/bad partition.  Thus
\begin{equation}
 tb(m+\ell)\le\frac{\mu P}{2},
 \qquad 0<t\le\frac{\mu}{100}\le\frac1{100}.
 \label{loc:selected-threshold}
\end{equation}
Every nonshort bad center has $V_i>tb$.  By
\eqref{loc:scale-condition} and \eqref{loc:V-bound},
\[
 m\le q_0+\frac{\sum_iV_i}{tb}
 \le\left(10+\frac6t\right)\frac Pb
 \le C_\mu\frac Pb.
\]
This proves the count bound.

\paragraph{Proxy geometry and the number of crowded balls.}
For each good $i$, choose a nearest facility $\phi_i\in S$ to $i$, and
write $r_i=d(i,\phi_i)$.  Such a center has $n_i>0$.  For every
$p\in C_i$, nearest choice and the triangle inequality give
$r_i\le d(i,s(p))\le d_p^*+s_p$.  Hence
\begin{equation}
 n_ir_i\le V_i\le tb,
 \qquad r_i\le t\operatorname{sep}_i.
 \label{loc:proxy-distance}
\end{equation}
Consider the open balls in the ambient metric
\[
 B_i=\{x:d(i,x)<\operatorname{sep}_i/4\},\qquad i\in O_g.
\]
They contain their respective proxies.  They are pairwise disjoint:
for distinct $i,j\in O$, both separations are at most $d(i,j)$, so
the sum of the ball radii is at most $d(i,j)/2$; membership in both
balls would contradict the triangle inequality.  In particular the
proxies are distinct.

There are $h-m$ good balls and $|S|=h+\ell$ local facilities.  Every good
ball contains at least one local facility.  Because the balls are
disjoint, the number of them containing at least two local facilities
is at most
\[
 |S|-(h-m)=m+\ell.
\]
Call these balls crowded.  All other good balls contain exactly their
one proxy as a local facility.

\paragraph{Charging the assignment excess.}
We upper-bound the sum of $d(p,\phi_i)-d_p^*$ over all good clusters
using three contributions.  Individual excesses may be negative;
only their total upper bound is needed.

First, for a good center with a crowded ball, the triangle inequality
and \eqref{loc:proxy-distance} give
\[
 \sum_{p\in C_i}\bigl(d(p,\phi_i)-d_p^*\bigr)
 \le n_ir_i\le tb.
\]
The total contribution from crowded balls is therefore at most
$tb(m+\ell)$.

Second, suppose $B_i$ contains exactly one local facility and
$p\in C_i$ is assigned in $S$ to a facility other than $\phi_i$.
That facility lies outside $B_i$, including possibly on its boundary,
and therefore
\[
 d_p^*+s_p\ge d(i,s(p))\ge\operatorname{sep}_i/4.
\]
The excess of this client is at most
\[
 d(p,\phi_i)-d_p^*\le r_i
 \le4t(d_p^*+s_p).
\]
Summing over these clients gives a total at most $4t\sum_iV_i$.

Third, consider the remaining clients of the singleton balls, namely
those assigned to their own proxy.  For such a good center $i$, put
\[
 A_i=\{p\in C_i:s(p)=\phi_i\},\qquad
 I_i=\{p\in D\setminus C_i:s(p)=\phi_i\}.
\]
We claim that
\begin{equation}
 \sum_{p\in A_i}(s_p-d_p^*)\le r_i|I_i|.
 \label{loc:incoming-swap}
\end{equation}
If $i=\phi_i$, the left side is zero.  Otherwise $i\notin S$, since a
nearest facility to a facility already in $S$ is that facility itself.
Exchange $\phi_i$ for $i$.  Send all clients previously assigned to
$\phi_i$ to $i$ and keep all other assignments.  Clients in $A_i$ have
change $d_p^*-s_p$; each incoming client in $I_i$ has change at most
$r_i$.  Local optimality implies
$0\le\sum_{p\in A_i}(d_p^*-s_p)+r_i|I_i|$, proving
\eqref{loc:incoming-swap}.  This also explains why overlap between
$O$ and $S$ creates no exception to the argument.

For $p\in I_i$, let $j\ne i$ be its assigned comparison center.
The path from $i$ to $j$ through $\phi_i$ and $p$ gives
\[
 d_p^*+s_p\ge d(i,j)-r_i
 \ge(1-t)\operatorname{sep}_i.
\]
Together with \eqref{loc:proxy-distance}, this bounds its incoming
charge by
\[
 r_i\le\frac{t}{1-t}(d_p^*+s_p).
\]
The proxies are distinct and the local assignment is fixed.  Thus a
client is incoming at at most one charged proxy, and summing
\eqref{loc:incoming-swap} gives a total at most
$\frac{t}{1-t}\sum_iV_i$ for the third contribution.

Combining the three contributions proves
\begin{align*}
 D_g-P_g
 &\le tb(m+\ell)+\left(4t+\frac{t}{1-t}\right)\sum_iV_i\\
 &\le\frac{\mu P}{2}
       +\left(24t+\frac{6t}{1-t}\right)P.
\end{align*}
Since $t\le1/100<1/2$, the last parenthesis is at most
$36t\le(9/25)\mu<\mu/2$.  This proves $D_g\le P_g+\mu P$ and finishes
the proof.  Empty comparison clusters were counted among the bad
centers throughout and never required division by $n_i$.
\end{proof}

For the main approximation theorem we henceforth use the specialization
$\mu=1/100$ and the absolute constant
\begin{equation}
 C_0:=C_{1/100}=10+60000\cdot2^{3999}.
 \label{loc:main-count-constant}
\end{equation}
Thus $C_0$ is fixed before any choice of the later algorithm parameters
$v$ and $L$.  Theorem~\ref{loc:anchoring} gives
$D_g\le P_g+0.01P$ and $m\le C_0P/b$ with this one choice.  The same
proof supplies smaller fixed proxy error in the refined recovery
theorem by choosing a different fixed $\mu$.

Two consequences will be useful.  If $|S|=h$, the number of facilities
of $S$ outside the distinct proxy set is exactly $m$, including when
some bad comparison clusters are empty.  Also, if the main
specialization satisfies $P_b\le P/10$, then
Lemma~\ref{loc:swap} gives
\begin{equation}
 \operatorname{cost}(S)
 \le\frac65(P-P_b+0.01P)+\frac{17}{5}P_b
 \le1.432P<1.6P.
 \label{loc:small-bad-cost}
\end{equation}
Finally, when all centers are declared bad, the proxy condition is
vacuous and $m=h$; that use of Lemma~\ref{loc:swap} needs no separation
assumption and applies also to $h=1$.
